\section{Экспериментальный план и постановка замеров}

\DTLloadrawdb{nlpparams}{tables/nlp_illustrative_params.csv}%
\newcommand{\NlpTabRows}{}%
\DTLforeach*{nlpparams}{\prm=param,\vl=colvalue}{%
  \eappto\NlpTabRows{\expandonce{\prm} & \expandonce{\vl}\\}%
}%

\subsection{Языковые данные (иллюстративная конфигурация)}

Основная серия экспериментов на естественном языке в полном объеме в данной работе не воспроизводится: ниже приведена \emph{типовая} постановка на публично доступных корпусах, чтобы зафиксировать масштабы и метрики. Численные значения в табл.~\ref{tab:main_params} являются \textbf{плейсхолдерами} для ориентира по порядку величин и не должны трактоваться как результат конкретного запуска из репозитория.

\begin{table}[H]
\centering
\caption{Иллюстративные параметры языкового эксперимента (публичные корпуса; числа условные; данные: \texttt{tables/nlp\_illustrative\_params.csv}; корпус WikiText-103 --- \cite{merity2016})}
\label{tab:main_params}
\begin{tabular}{l p{0.72\linewidth}}
\toprule
Параметр & Значение (иллюстрация) \\
\midrule
\NlpTabRows
\bottomrule
\end{tabular}
\end{table}

Для такой постановки на валидации ожидаются, например, перплексия порядка \(28\)--\(32\) и доля точного совпадения верхнего токена порядка \(0{,}35\)--\(0{,}42\) при фиксированной глубине рекурсии; эти диапазоны приведены только как ориентир методологии оценки и подлежат замене фактическими значениями при подключении языкового трека к тому же коду, что использован для ARC-AGI.

\subsection{Сравниваемые варианты}

Сравнение должно быть устроено так, чтобы не смешивать вклад разных частей. Минимальный набор вариантов приведен в табл.~\ref{tab:model_variants}. Первая строка задает обычную модель с фиксированной глубиной. Остальные строки проверяют, что именно дает выигрыш: предобученная основа, рекурсивное уточнение, модель распределения времени или правило остановки.

\begin{table}[H]
\centering
\caption{Сравниваемые варианты модели}
\label{tab:model_variants}
\begin{tabularx}{\textwidth}{lYYY}
\toprule
Вариант & Основа & Рекурсия & Время остановки \\
\midrule
Фиксированная глубина & малая языковая модель & нет & не используется \\
Полный запуск рекурсии & малая языковая модель & до \(K\) шагов & всегда \(K\) \\
Ранняя остановка по энтропии & малая языковая модель & до \(K\) шагов & простая эвристика \\
Распределение \(\tau^\star\) & малая языковая модель & до \(K\) шагов & одно из пяти правил \\
Низкоранговая донастройка & дообученная основа & нет или есть & отдельно сравнивается \\
\bottomrule
\end{tabularx}
\end{table}

\subsection{Исследование вклада компонентов}

Исследование вклада компонентов нужно ограничить только тем, что относится к основной гипотезе. Проверяются четыре фактора: форма распределения времени, тип разметки, правило остановки и способ переиспользования малой языковой модели. Шаблон таблицы приведен в табл.~\ref{tab:ablation_template}.

\begin{table}[H]
\centering
\caption{Шаблон таблицы исследования вклада компонентов}
\label{tab:ablation_template}
\begin{tabularx}{\textwidth}{lYYYY}
\toprule
Вариант & Потеря речи & \(\mathrm{NLL}_{\tau}\) & Среднее число шагов & Экономия \\
\midrule
Конечное распределение + накопленная вероятность & \placeholder & \placeholder & \placeholder & \placeholder \\
Сглаженная разметка + накопленная вероятность & \placeholder & \placeholder & \placeholder & \placeholder \\
Смесь геометрических + малое будущее улучшение & \placeholder & \placeholder & \placeholder & \placeholder \\
Отрицательное биномиальное + квантиль & \placeholder & \placeholder & \placeholder & \placeholder \\
Бюджетное правило при заданном бюджете & \placeholder & \placeholder & \placeholder & \placeholder \\
\bottomrule
\end{tabularx}
\end{table}

\subsection{Реализованный пайплайн ARC-AGI}

Для ARC-AGI в репозитории зафиксирован свип по шести семействам распределения времени остановки (конфигурации \texttt{configs/oracle\_sweep\_arc\_agi/}). Для каждого семейства обучается рекурсивная модель с головой оракула; затем на валидации запускается сетка правил остановки (скрипт \texttt{validate\_oracle\_sweep.sh}): стратегии \texttt{cumulative\_probability}, \texttt{future\_improvement}, \texttt{hazard}, \texttt{quantile}, \texttt{budget} с порогами и бюджетами из скрипта, политики ответа \texttt{last} и \texttt{argmax\_interval}. Итог агрегируется в \texttt{stopping\_summary.csv} и \texttt{stopping\_summary.md} (метрика отбора лучшей строки сетки: \texttt{token\_acc\_best\_of\_policies}, режим \texttt{max}).

Дополнительно для каждого базового прогона выполняется дообучение \emph{только} головы оракула (\texttt{finetune\_oracle\_sweep.sh}, подкаталог \texttt{oracle\_finetune/}), после чего оценка останавливающих правил повторяется на том же валидационном сплите.

Численные результаты свипа и сравнение до/после дообучения приведены в разд.~\ref{sec:arc_results}; эмпирическое распределение оракульного \(\tau^\star\) на валидации иллюстрируется гистограммой для конфигурации \texttt{finite\_discrete}.

